\subsection{The equality boundary in characteristic two}
\label{subsec:positive-excess-boundary}

We now prove that the first positive-excess bound can be attained
in characteristic two only under a further restriction on the curve.
The proof examines equality in the resolution estimate. The resolved
curve has square zero, so a conormal coefficient need not vanish.
Its descent, together with a leading two-form on the original
surface, instead detects a nontrivial two-torsion line bundle on
the curve.

\begin{proof}[Proof of the discrepancy and nonvanishing assertions in Corollary~\ref{cor:positive-excess-characteristic-two-boundary}]
Write $C=Y$, $N=\OO_C(C)$ and $g=g(C)$. Thus
\[
 2g-2=(2a+1)d,\qquad
 \mathcal T=\omega_C\otimes N^{-(2a+1)},\qquad
 \deg\mathcal T=0.
\]
The lifting argument in the proof of
Theorem~\ref{thm:positive-excess-picard} gives a line bundle $L$
with $L^{\otimes2}\simeq M$. It has exact order four.
Resolve $X$ away from $C$, as in that proof. All the displayed
curve data remain unchanged, so we may assume that $X$ is smooth.

\smallskip\noindent
\emph{A finite model near the curve.}
Use frames of $L$ and their squares as frames of $M$. If $g_{ij}$
are the transitions of these frames of $M$, then $dg_{ij}=0$.
Let $f_i$ be the units defining the chosen trivialization of
$M^{\otimes2}$, and let $t_i=0$ define $C$. The local
$p$-basis calculation in
Lemma~\ref{lem:positive-square-root-genus} gives
\begin{equation}\label{eq:boundary-square-approximation}
 f_i=b_i^2+t_i^{2a}A_i,
\end{equation}
where $b_i$ is a unit after shrinking around $C$.
Since $f_i=g_{ij}^2f_j$, we have
$b_i-g_{ij}b_j\in I_C^a$.
Consequently the finite flat double covers
\[
 y_i^2=A_i,\qquad z_i=b_i+t_i^ay_i
\]
glue by
\[
 y_i=g_{ij}(t_j/t_i)^a y_j+
       \frac{g_{ij}b_j-b_i}{t_i^a}.
\]
Denote the resulting surface near $C$ by $V$ and its curve over $C$
by $B_V$. The tangential leading coefficient of the logarithmic
form of $M$ is nonzero, so $d(A_i|_C)\ne0$ generically.
Thus $B_V$ is generically a field extension of degree two.
Its algebra is locally free over the integral curve $C$, hence
torsion-free; generic integrality makes $B_V$ integral.

The submodule of constants in that rank-two algebra has quotient
with transition functions $g_{ij}(t_j/t_i)^a$. Since $M|_C$
is trivial, this gives the exact sequence of $\OO_C$-modules
\[
 0\longrightarrow\OO_C\longrightarrow
 (\pi_V)_*\OO_{B_V}\longrightarrow N^a\longrightarrow0.
\]
Riemann--Roch on $C$ therefore gives
\begin{equation}\label{eq:boundary-finite-genus}
 p_a(B_V)-g=g-1-ad=\frac d2.
\end{equation}
In particular $d$ is even.

Let $\pi:W\to X$ be the normalization of the original torsor.
Near $C$ it is also the normalization of $V$.
By Lemma~\ref{lem:positive-square-root-genus},
$B=\pi^*C$ is reduced, integral and Cartier, $B^2=2d$,
and its normalization has genus $g$.
Apply the resolution procedure in
Lemma~\ref{lem:positive-square-resolution-loss}: first take a
minimal ambient resolution $S_0\to W$, and then blow up singular
points of the strict curve until its transform $D$ on a smooth
projective surface $S$ is smooth. The map $q:D\to C$ is finite,
purely inseparable of degree two, and is relative Frobenius
\cite[Tag~0CCZ]{Stacks}. The preceding genus bound gives
\begin{equation}\label{eq:boundary-nonnegative-square}
 D^2\geq2d-2\delta=0.
\end{equation}

\smallskip\noindent
\emph{The conormal leading coefficient.}
The pullback $L'$ of $L$ to $S$ has square trivialized by the
units $z_i$. Its logarithmic form is the global regular form
$\eta=dz_i/z_i$. At the generic point of $C$, omit indices.
The differentials $dt,dA$ form a basis over $k(X)$, so write
\[
 db=R\,dt+S_1\,dA,\qquad R,S_1\in k(X).
\]
On the cover $A=y^2$, and $a$ is odd. Hence
\begin{equation}\label{eq:boundary-root-differential}
 \eta=\frac{(R+t^{a-1}y)\,dt+t^a\,dy}{z}.
\end{equation}
Here $t$ is a uniformizer along $D$, and $dt,dy$ form a
cotangent basis at its generic point. In particular $\eta\ne0$.
A trivial line bundle on a proper integral variety has zero
logarithmic differential for a square trivialization, because its
global units are constants. Thus $L'$ has exact order two.
The image of $\eta$ in $\Omega_D^1$ is zero: $dt$ is conormal
and the $dy$ coefficient is divisible by $t^a$.
The injectivity of the logarithmic-form map on torsion line bundles
of a smooth proper curve, established in the proof of
Lemma~\ref{lem:positive-square-resolved-transfer}, shows that
$L'|_D$ is trivial. None of these assertions requires $D^2>0$.

The residue $y_0$ lies in $k(D)=k(C)^{1/2}$ and not in $k(C)$.
Every normalized leading residue of $R$ lies in $k(C)$.
Therefore the two terms of $R+t^{a-1}y$ cannot cancel when
their orders coincide. If $h=\operatorname{ord}_D(\eta)$, then
\begin{equation}\label{eq:boundary-conormal-order}
 0\leq h=\min\{v_C(R),a-1\}\leq a-1,
\end{equation}
where the valuation of zero is infinity. Regularity gives the
lower bound. The leading coefficient is purely conormal, and
defines a nonzero global section
\begin{equation}\label{eq:boundary-conormal-section}
 \sigma\in H^0\bigl(D,N_{D/S}^{-(h+1)}\bigr).
\end{equation}
Here division by the $h$th power of a local equation is regular
along the entire smooth curve: in the regular local rings, this
divisibility can be checked at its generic point.
If $D^2>0$, this line bundle would have negative degree.
Combining with \eqref{eq:boundary-nonnegative-square} gives
$D^2=0$. The section $\sigma$ is therefore nowhere zero.

\smallskip\noindent
\emph{Equality in the genus and resolution losses.}
Let $B_0$ be the strict curve on $S_0$, and put
\[
 e=p_a(B_V)-p_a(B),\quad
 u=p_a(B)-p_a(B_0),\quad
 v=p_a(B_0)-g.
\]
The first number is the length of the quotient for the finite
birational map $B\to B_V$, and the other two are the
nonnegative genus losses in the resolution proof.
By \eqref{eq:boundary-finite-genus}, $e+u+v=d/2$.
Writing the ambient pullback as $B_0+F$,
\eqref{eq:positive-square-genus-loss} gives
$B^2-B_0^2=-F^2\leq2u$.
If the later blowup multiplicities are $\mu_j\geq2$, then
\[
 v=\sum_j\frac{\mu_j(\mu_j-1)}2,
 \qquad B_0^2-D^2=\sum_j\mu_j^2\leq4v.
\]
It follows that
\[
 0=D^2\geq2d-2u-4v=4e+2u.
\]
Thus $e=u=0$, $F^2=0$, and every $\mu_j$ equals two.

The surface $W$ is smooth along $B$. Indeed, if a singular
point lay on $B$, its local Cartier equation would have positive
order on every exceptional component over that point. Properness
forces $B_0$ to meet the exceptional fibre, so $B_0\cdot F>0$.
But $(B_0+F)\cdot F=0$ would give $F^2=-B_0\cdot F<0$,
a contradiction.
Since $e=0$, the finite birational map $B\to B_V$ is an
isomorphism. To see this directly, the quotient of the structure
sheaves has finite length; both integral proper curves have global
functions $k$, so equality of arithmetic genera makes that length
zero \cite[Tag~0CE4]{Stacks}.
The normalization $W\to V$ is then an isomorphism near $B_V$.
For a finite normalization extension $A_0\subset A_1$ near the
curve, the isomorphism modulo its equation $t$ gives
$A_1=A_0+tA_1$. Nakayama applied to $A_1/A_0$ proves the
claim. In particular $V$ is smooth along $B_V$.

\smallskip\noindent
\emph{The normal line after the cusps are resolved.}
Let $\alpha$ be the tangential leading coefficient of the
logarithmic form of $M$, and put $Z=\operatorname{div}(\alpha)$.
Then
\begin{equation}\label{eq:boundary-leading-divisor}
 \OO_C(Z)=\omega_C\otimes N^{-2a},\qquad \deg Z=d>0.
\end{equation}
At a point $P$ of $C$, the completed curve equation is
$y^2=A(x,0)$. The multiplicity of $Z$ is the order of
$d(A(x,0))$. In characteristic two that order is even, say
$2\ell$. The even powers of $A(x,0)$ form a square, while
the odd part is $x^{2\ell+1}$ times a square unit. Subtracting
the first square from $y$ and dividing by the unit gives
\[
 w^2=x^{2\ell+1}.
\]
For $\ell=0$ this is smooth. For $\ell\geq1$ its normalization
is $x=s^2$, $w=s^{2\ell+1}$, and its genus defect is $\ell$:
the missing exponents of $k[[s^2,s^{2\ell+1}]]$ are
$1,3,\ldots,2\ell-1$.

Because $V$ is smooth along the curve, at a cusp $x,w$ are
surface parameters. Exactly $\ell$ point blowups resolve the
curve. At each step its multiplicity is two and the exceptional
equation restricts to $s^2$ on the normalization. Thus
$D_{\rm new}=b^*D_{\rm old}-2E$ subtracts $4P'$ from the
normal line at each step, where $P'$ lies above $P$.
The point-blowup normal bundle is as in
\cite[Tag~0AGQ]{Stacks}; the coefficient here also follows
directly by substituting $w=xw_1$ successively.
The total correction is $4\ell P'$, exactly the coefficient
of $q^*Z$, since $q$ has ramification index two. Therefore
\begin{equation}\label{eq:boundary-normal-line}
 \begin{aligned}
 N_{D/S}
 &\simeq q^*N\otimes\OO_D(-q^*Z)\\
 &\simeq q^*\bigl(N\otimes\OO_C(-Z)\bigr)
 \simeq q^*(\mathcal T^{-1}).
 \end{aligned}
\end{equation}
Away from $Z$ this is the natural identification before any
blowup correction.

\smallskip\noindent
\emph{A cusp forces conormal order zero.}
The divisor $Z$ has positive degree, so choose one of the cusps above.
The form $\eta=dz/z$ is regular on $V$ before the cusp resolution.
At a cusp with complete equation $f=w^2-x^{2\ell+1}$, write
\[
 \eta=f^h\bigl(A(x,w)\,dx+B(x,w)\,dw\bigr).
\]
The coefficients $A,B$ are regular power series: divisibility by
the power of the prime Cartier equation can be checked at its
generic point in the regular local ring. On the branch
$x=s^2$, $w=s^{2\ell+1}$, the purely conormal leading condition
gives $s^{2\ell}B(s^2,s^{2\ell+1})=0$. Thus $B=fB_1$.
The composite chart of the $\ell$ cusp blowups is
\[
 w=x^\ell v,\qquad H=x-v^2,\qquad
 f=-x^{2\ell}H,\qquad dx=dH.
\]
After division by $H^h$ and restriction to the strict curve,
the term involving $fB_1$ vanishes. The conormal coefficient is
\[
 v^{4\ell h}A(v^2,v^{2\ell+1}).
\]
It would vanish at $v=0$ if $h>0$, contradicting the nowhere-zero
section $\sigma$. Hence $h=0$. Completion preserves the branch
normalization and finite-length defect
\cite[Tags~0C1R and~0C3V]{Stacks}; faithful flatness detects this zero
on the original curve.

\smallskip\noindent
\emph{For $a\geq3$ the discrepancy line is trivial.}
Suppose $a\geq3$. Then $h=0<a-1$. On the original surface put
$\theta_i=d\log b_i$. The overlap relation and $dg_{ij}=0$
give
\[
 \theta_i-\theta_j\in I_C^{a-1}\Omega_X^1.
\]
Write $\omega=d\log M$. Since
$\omega\in I_C^{2a}\Omega_X^1$, the two-forms
$\theta_i\wedge\omega$ agree modulo
$I_C^{3a-1}\Omega_X^2$.
At the generic point, \eqref{eq:boundary-square-approximation}
gives
\[
 \theta\wedge\omega
   =\frac{t^{2a}R}{bf}\,dt\wedge dA.
\]
Its exact order is $2a+h<3a-1$. Its leading coefficient
therefore glues to a nonzero global section of
\[
 \omega_X|_C\otimes N^{-(2a+h)}
 \simeq\omega_C\otimes N^{-(2a+h+1)}
 =\mathcal T\otimes N^{-h}.
\]
This line has degree $-hd$. Hence $h=0$, and its nonzero
section trivializes $\mathcal T$. Since $h<a-1$ and $a$ is
odd, this case requires $a\geq3$.

\smallskip\noindent
\emph{For $a=1$ the discrepancy line has exact order two.}
Suppose $a=1$, so $h=0=a-1$. Equations
\eqref{eq:boundary-conormal-section} and
\eqref{eq:boundary-normal-line} show that $\sigma$ trivializes
$q^*(\mathcal T^a)$. The norm along the finite flat
degree-two curve map gives $\mathcal T^{2a}\simeq\OO_C$.
We must also show that $\mathcal T^a$ is nontrivial.

In the conormal frame arising from $t$, the scalar of $\sigma$
at the generic point is
\[
 \gamma=\frac{c+y_0}{b_0},\qquad
 c\in k(C),\quad b_0\in k(C)^*,\quad y_0\notin k(C).
\]
Here $c$ is the residue of $t^{1-a}R$, or zero if that
element has positive valuation. Thus $\gamma\notin k(C)$.
By the natural compatibility of
\eqref{eq:boundary-normal-line} away from $Z$, replacing the
frame by one pulled back from a rational frame of $\mathcal T^a$
multiplies $\gamma$ only by an element of $k(C)^*$.
So $\sigma$ does not descend as a rational section on $C$.
If $\mathcal T^a$ were trivial, choose a nowhere-zero section
$\tau$ on $C$. The ratio $\sigma/q^*\tau$ would be a global
regular function on the integral projective curve $D$, hence a
constant in $k$. This contradicts non-descent. Consequently
$\mathcal T^a$ has exact order two.
Since $a=1$, the line $\mathcal T$ itself has exact order two.

\smallskip\noindent
\emph{Two-torsion on the curve in the remaining case.}
When $\mathcal T^a$ is nontrivial it already supplies the required
two-torsion. Suppose instead that $h=0$ and $\mathcal T$ is
trivial. The functions $b_i|_C$ trivialize $(L|_C)^{\otimes2}$;
therefore $\beta_i=d\log(b_i|_C)$ glue to its logarithmic form
$\beta$. The zero divisor $Z$ is nonempty by
\eqref{eq:boundary-leading-divisor}. At a point of $Z$, choose
coordinates $x,t$. There $A_x=0$. Smoothness of
$y^2=A(x,t)$ forces $A_t\ne0$. The nowhere-zero leading
coefficient of $\theta\wedge\omega$ is, up to a unit,
\[
 b_xA_t-b_tA_x.
\]
At this point it equals $b_xA_t$, so $b_x\ne0$. Thus
$\beta$ is nonzero, and $L|_C$ is a nontrivial line bundle
of exact order two. This proves $\Pic(C)[2]\ne0$ in both
cases. Combining with Theorem~\ref{thm:positive-excess-picard}
gives $\delta>d$ whenever $\Pic(C)[2]=0$.
The stated integer bounds follow because $\delta$ is even.
\end{proof}

The smallest case is $m=2$: equality would require the specific
line $\omega_C\otimes\OO_C(-3C)$ to have exact order two.
For $m=6$, the discrepancy line must be trivial. Neither calculation constructs a surface with these
properties. The rank-two refinement at $m=2$ is proved in
Proposition~\ref{prop:smallest-endpoint-rank}; its strictness consequence
completes Corollary~\ref{cor:positive-excess-characteristic-two-boundary}.
The actual conormal trivialization above, together with the
positive-characteristic moving lemma, yields the fibration in
Proposition~\ref{prop:boundary-resolved-fibration}.
